% Generated draft: requires mathematical and source review.
\documentclass[11pt]{article}
\usepackage[margin=25mm]{geometry}
\usepackage{fontspec}
\setmainfont{Latin Modern Roman}
\usepackage{amsmath,amssymb,mathtools,unicode-math}
\setmathfont{Latin Modern Math}
\usepackage{enumitem,needspace,xurl,hyperref}
\setlength{\parindent}{0pt}
\setlength{\parskip}{5pt}
\setlength{\emergencystretch}{3em}
\newcommand{\N}{\mathbb{N}}
\newcommand{\Z}{\mathbb{Z}}
\newcommand{\Q}{\mathbb{Q}}
\newcommand{\R}{\mathbb{R}}
\newcommand{\C}{\mathbb{C}}
\newcommand{\F}{\mathbb{F}}
\newcommand{\A}{\mathbb{A}}
\newcommand{\PP}{\mathbb{P}}
\newcommand{\E}{\mathbb{E}}
\newcommand{\eps}{\varepsilon}
\newcommand{\dd}{\,\mathrm{d}}
\begin{document}
\section*{UnsolvedMath ID: 20002321}
\textbf{Source problem code:} catalogue entry
\subsection{Definitions and mathematical statement}
Question 19 asks what additional information is needed to build a computable prime model of a decidable theory \(T\) with countably many countable models, noting that \(\omega\)-stability suffices. The count is understood up to isomorphism in the aid's usual countable-language setting. The source does not explicitly settle completeness or whether ``countably many'' includes finitely many. The following DCF/Turing-degree heading belongs to the next section and does not restrict this question or attribute it to Calvert.

For the aid's conditional result assume \(T\) is a complete decidable atomic theory in a computable countable language. A complete type is principal when isolated by a consistent formula; atomic means principal types are dense among consistent formulas. A prime model admits an elementary embedding into every model of \(T\). A computable copy has decidable atomic diagram, while a decidable copy has decidable full elementary diagram.

Assume further that the prime-model isomorphism type is relatively decidable: for every copy \(P\) of that type, its elementary diagram is computable from its atomic diagram. The aid reports equivalence of:
\begin{enumerate}
\item a computable prime copy of \(T\) exists;
\item a decidable prime copy exists;
\item the complete principal types, across finite arities, have a uniform computable listing;
\item there is a computable principal-type selector returning indices for total computable characteristic functions.
\end{enumerate}
For clarity, a uniform type listing supplies total effective membership tests for the listed complete types and includes every principal type, allowing repetitions. The selector, in the aid's intended formulation, takes a consistent formula \(\varphi(\bar x)\) and returns an index for a complete principal type containing it, whose characteristic function on formulas is total. Merely listing isolating candidates or producing partial indices is not that conclusion. The summary does not reproduce its exact coding/uniformity definitions, so these details require checking against the full statement.

The aid further reports that every c.e. model-complete theory is uniformly relatively decidable: elementary diagrams are computable uniformly from atomic diagrams. Model-complete means every embedding between models is elementary. Consequently all four conditions are equivalent for every complete decidable atomic model-complete theory, including that portion of the few-countable-models setting. The extra relative-decidability hypothesis is not asserted automatic for every complete decidable atomic theory, and the equivalence alone does not say its four conditions always hold. The unrestricted workshop question remains open in the supplied assessment.
\subsection{Short English statement}
The source asks how to obtain computable prime models in the few-model setting. The aid gives equivalent effective type-selection conditions when the prime type is relatively decidable, with model-completeness supplying that extra hypothesis. It distinguishes computable atomic diagrams from decidable elementary diagrams and does not claim a general construction for the original class.
\subsection{Sources}
\begin{itemize}
\item Local UnsolvedMath registry, permanent ID 20002321, exact source code catalogue entry; canonical \texttt{aim-logic-notes.json}, \texttt{notes[96]}, tag \texttt{question}; AIM Workshop Problem Lists, Logic, \emph{Computable stability theory}, Question 19. Created and updated 14 August 2026.
\item Source URL \url{https://aimath.org/pastworkshops/computestabproblems.pdf}. The supplied verification separates the following ``Turing degrees of DCFs (Calvert)'' heading. The original notes \(\omega\)-stability as sufficient. The PDF was not accessed.
\item Attempt 1: \texttt{conditional\_result}, problem status at run \texttt{open}, rigor label \texttt{proved}, novelty confidence \texttt{low}; classification \texttt{PARTIAL-PROGRESS}, suggested difficulty L4, registry status \texttt{partially\_solved}. Candidate: principal-type-selector equivalence under relative decidability, including the model-complete specialization. No full authoritative theorem references are supplied.
\item Ulam AI / UnsolvedMath Contributors; CC BY 4.0. Machine-generated research aid, not peer reviewed; independently verify before citation.
\end{itemize}
\end{document}
